\chapter{Discussion \& Limitations}
\label{ch:discussion}

%==============================================================================
\section{Inductive Bias Against Capacity in Sample-Constrained Neuroimaging}
\label{sec:disc-capacity}

The clearest pattern across this work is that, in this regime, capacity is a liability and
constraint is a feature.

The evidence is consistent across three independent comparisons. A transformer attending
over ${\sim}1{,}381$ tokens per subject, trained on ${\sim}76$ subjects per fold, achieves
strong validation performance that does not transfer --- a CV-to-test gap of $-0.198$,
against $+0.060$ for the leanest recurrent model (Table~\ref{tab:capacity}). Training a
graph encoder end-to-end from random initialisation on the same data produces the largest
fold-to-fold variance of any arm and, in at least two instances, outright optimisation
collapse. And freezing a pretrained encoder outperforms fine-tuning it, which is to say that
\emph{removing} degrees of freedom improved the result.

This is not a claim that transformers or graph encoders are bad models. It is a claim about
a specific and very common regime: when the number of labelled subjects is in the low
hundreds and the feature dimension is in the tens of thousands, the model class that wins is
the one whose hypothesis space the data can actually constrain. The literature's move toward
foundation models pretrained on $10^4$--$10^5$ subjects (Section~\ref{sec:longitudinal}) is
the other coherent response to the same problem --- borrow the constraint from elsewhere ---
and is not available to a single-site cohort.

%==============================================================================
\section{What Reconstruction Pretraining Actually Bought}
\label{sec:disc-pretraining}

The premise of Stage~1 was that a self-supervised reconstruction objective would learn a
representation transferring to conversion prediction. The honest reading of the ablation is
narrower: pretraining's measurable benefit is \textbf{optimisation stability, not peak
performance}.

Pretrained initialisation clearly beats random initialisation when the encoder is trained
end-to-end --- the random arm is the worst and most variable of the four. But a model with
\emph{no encoder at all}, feeding mean-pooled raw connectivity straight to the recurrent
core with roughly seventeen times fewer trainable parameters, is not beaten by the
pretrained frozen pipeline. The pretraining stabilises a problem that the no-encoder
configuration never has.

The mechanism proposed in Section~\ref{sec:graph-construction} accounts for this without
appeal to novelty. The node features already contain the full connectivity profile at full
precision. Routing them through $k$NN-8 sparsification, binarisation, and a 64-dimensional
bottleneck optimised for edge reconstruction is a sequence of lossy steps applied to
information that was already present --- and the reconstruction objective supplies no
gradient preserving the between-subject variance the classifier needs
(Section~\ref{sec:gap-pretext}).

This aligns with, and extends, the strongest recent evidence in the field: the finding that
message aggregation ``consistently degrades'' predictive performance across four large-scale
cross-sectional neuroimaging studies (Section~\ref{sec:gap-aggregation}). The contribution
here is not the direction of that result --- it is that the result now also holds in the
longitudinal, small-$N$ clinical regime that study did not cover.

%==============================================================================
\section{On ``Stability'' Claims About Non-Deterministic Baselines}
\label{sec:disc-determinism}

The methodological finding of Chapter~\ref{ch:sota-repair} generalises well beyond
Brain-TokenGT, and is the part of this thesis most likely to be useful to other people.

An intervention that stops a model from crashing is not the same as an intervention that
makes it stable, and the two are easy to conflate because they share the same evidence:
runs that complete. This project made exactly that conflation --- labelling an arm
``stabilised'' and reporting a four-seed standard deviation for it --- before an audit of
its own run artifacts showed within-seed variance exceeding between-seed variance across 19
runs.

The general lesson is a one-line addition to any evaluation protocol: \textbf{before
reporting a multi-seed standard deviation, run one seed several times.} If the within-seed
spread is comparable to or larger than the between-seed spread, the reported error bar does
not describe seed sensitivity, and no significance test built on it is valid. This costs a
handful of extra runs and is, as far as this review could establish, absent from the
brain-graph literature entirely.

It cuts both ways, and this thesis is not exempt: the frozen-encoder arm used in the
matched-cohort comparison has exactly one run per seed, so its own within-seed variance is
argued rather than measured (see the exposure noted in
Section~\ref{sec:btgt-reproducibility}).

%==============================================================================
\section{Threats to Validity}
\label{sec:disc-threats}

\paragraph{Cohort size dominates everything.} With 34 held-out test subjects, of whom 14 are
converters, one subject changing side of the boundary is worth roughly $0.03$--$0.04$ AUC.
Most of the differences discussed in this thesis are of that magnitude. This is why the
argument rests on cross-validation distributions, effect-size ranges, and consistency across
arms rather than on test point estimates --- and why the ADNI replication is load-bearing
rather than decorative.

\paragraph{The label is a censored observation in disguise.} A subject labelled ``stable''
may simply not have converted yet. The binary framing therefore mixes true non-converters
with not-yet-converters in proportions that depend on follow-up length --- which, in this
cohort, differs systematically between the classes.

\paragraph{Informative attrition.} Converters are followed longer
(Section~\ref{sec:gap-attrition}). The within-subject slope analysis argues that the models
accumulate evidence rather than counting visits, but that analysis is a refutation of one
specific shortcut, not a proof that no length-related confound remains.

\paragraph{Single-site development.} Every architectural and hyperparameter decision was
made against DELCODE. Even with an untouched test split, the accumulation of design choices
across a project is a form of soft overfitting to a cohort that no held-out split detects.
Only external validation addresses it.

\paragraph{A negative result is not proof of absence.} The finding that the graph encoder
does not earn its place is a statement about \emph{this} encoder, \emph{this} graph
construction, and \emph{this} cohort size. A construction retaining edge weights, or a
pretext objective aligned with discrimination, might well behave differently. That is a
hypothesis this work motivates rather than one it tests.

%==============================================================================
\section{Clinical Relevance \& Translation}
\label{sec:disc-clinical}

Two gaps separate these results from clinical utility, and neither is closed by a better
AUC.

First, the output is a binary flag over an undefined horizon, not a calibrated time to
event. A clinician deciding whether to start an anti-amyloid therapy needs a risk over a
stated interval; ``will convert by end of follow-up'' is not that, and the survival
extension that would supply it is scoped rather than delivered
(Section~\ref{sec:prognoser-scope}).

Second, the operating point matters more than the ranking. In prodromal screening a false
negative costs a patient an intervention window, while a false positive costs distress and
a downstream PET or lumbar puncture. Which error to prefer is a clinical judgement, not a
modelling one, which is why thresholds here are reported explicitly and derived from
validation only --- and why probability calibration, not just discrimination, is reported.

Third and most practically: nothing in this thesis has been evaluated prospectively, on a
new scanner, or against a clinician's own judgement. The comparison that would matter ---
does this add information beyond what a memory-clinic physician already has from cognitive
testing, age, and APOE status --- has not been run. The metadata-only floor
(CV AUC $0.6157$) is the nearest available proxy, and it is a weak one.
